\documentclass[10pt,a4paper]{amsart}
\usepackage[margin=19mm]{geometry}
\usepackage{amsmath,amssymb,mathtools}
\usepackage[scr=rsfs,cal=euler]{mathalfa}
\usepackage{xfrac}
\usepackage[unicode,hidelinks,bookmarks=false]{hyperref}
\newcommand{\ie}{i.e.\ }
\newcommand{\cf}{cf.\ }
\newcommand{\resp}{respectively\ }
\newcommand{\Subsection}{\subsection}
\providecommand{\nobreakdash}{\nobreak\hbox{-}}
\setlength{\emergencystretch}{2em}
\multlinegap0pt
\usepackage{amssymb}
\newtheorem{lemma}{Lemma}
\newcommand{\R}{R}
\newcommand{\T}{\mathcal T}
\title{Divergence of decreasing rearranged Fourier series\\ on every infinite compact abelian group}
\author{Morten Nielsen}
\date{Source: arXiv:2609.03506v1 (CC BY 4.0)}
\begin{document}
\maketitle

\noindent\emph{Source and licence.} Divergence of decreasing rearranged Fourier series\\ on every infinite compact abelian group. Version arXiv:2609.03506v1 (CC BY 4.0), distributed by arXiv under the Creative Commons Attribution 4.0 International licence, http://creativecommons.org/licenses/by/4.0/, as recorded in the arXiv licence metadata for this exact version and shown on the abstract page https://arxiv.org/abs/2609.03506v1. Full licence notice: \texttt{licenses/arxiv-2609.03506-CC-BY-4.0.txt}.\par\smallskip

\noindent\textit{Editorial scope.} Counted unit: the article's lemma lem:realize (source0.tex lines 371--398). The article's complete original proof of it is delivered with it (source0.tex lines 400--567, one \texttt{proof} environment of 168 lines). No mathematics is written, repaired or re-derived: the statement and the proof are the article's own lines, in the article's own notation, and no retained line is substituted or re-flowed.  Retained uncounted as the source-local context the proof needs: lines 290--304.  Nothing was omitted from any retained span, so the package carries no omission record.  No retained line contains a citation, so the wrapper carries no bibliography and no bibliography entry.  No retained line contains a figure environment, an image-inclusion command or drawing code, so no image is delivered and the package declares no asset dependency.  Editorial formatting only, in addition: an AMS \texttt{amsart} preamble that restates the article's own declarations and macros used by the retained lines, namely the environments \texttt{lemma} on separate counters, together with the standard packages those lines need.  The retained environments print in this document as Lemma 1 in order of appearance, whereas the article prints the same items with the numbering of its own full document.  The complete native source of the article as distributed in the arXiv e-print for this exact version is bundled unchanged as \texttt{sources/209302/source0.tex}.
\begin{lemma}\label{lem:chirp}
Let \(E\subseteq\R_p^n\) and \(d\geq n\). Write
\(X=(x,z)\in\R_p^n\times\R_p^{d-n}\), and let \(Y\in\R_p^d\).
The function
\[
 F_E(X,Y)=1_E(x)\omega^{X\cdot Y}
\]
satisfies \(|F_E|=1_E(x)\) and
\begin{equation}
 \widehat F_E(a,b)=
 p^{-d}1_E(b_1,\ldots,b_n)\omega^{-a\cdot b}.
\end{equation}
In particular, \(|\widehat F_E(a,b)|\in\{0,p^{-d}\}\), also when
\(E=\varnothing\).
\end{lemma}

\begin{lemma}\label{lem:realize}
Fix \(N\geq1\) and any bijection \(r:\T_N\to\{1,\ldots,L\}\).
There exist an integer \(D\geq1\), an integer \(C\),
sets \(E_v\subseteq\R_p^D\) for \(|v|\leq N\), a function
\(k:\R_p^D\to\{1,\omega,\ldots,\omega^{p-1}\}\), and functions
\(P_v:\R_p^D\to\mathbb C\) for \(v\in\T_N\), each a character polynomial, with the following properties:
\begin{enumerate}
 \item \(E_\varnothing=\R_p^D\). For every \(v\in\T_N\), the \(p\) sets
       \(E_{va}\), \(0\leq a\leq p-1\), partition \(E_v\). Moreover,
       \(\mu_D(E_v)=p^{-|v|}\) for every \(v\) with \(|v|\leq N\).
 \item For \(v\in\T_N\), the function
       \[
        h_v:\R_p^D\to\{0,\zeta_0,\ldots,\zeta_{p-1}\},\qquad
        h_v=\sum_{a=0}^{p-1}\zeta_a1_{E_{va}},
       \]
       vanishes outside \(E_v\) and equals \(\zeta_a\) on \(E_{va}\);
       and \(P_v=k\,h_v\).
 \item The Fourier supports of the \(P_v\) are pairwise disjoint.
 \item The exponents \(C+r(v)\) are nonnegative, and, for
       \(\xi\in\R_p^D\),
       \[
        |\widehat P_v(\xi)|\in\{0,p^{-(C+r(v))}\}.
       \]
       Moreover, \(\widehat P_v(0)=0\).
       Thus the nonzero coefficients have equal modulus within each
       node, while these moduli strictly decrease as \(r(v)\) increases.
\end{enumerate}
\end{lemma}

\begin{proof}
We proceed in rounds \(j=0,1,\ldots,N-1\). At the end of round \(j\),
the polynomials \(P_v\) have been constructed for all \(|v|\leq j\),
and the sets \(E_v\) have been constructed for all \(|v|\leq j+1\).
Thus the sets needed to build the depth-\((j+1)\) polynomials are
available when the next round begins. Previously constructed objects
are always lifted to an enlarged product group by ignoring the new
coordinates. Their current versions may be multiplied by common
chirps in later rounds; the sets \(E_v\) and the functions \(h_v\)
remain unchanged.

\emph{Round \(0\).}
On \(\R_p\), let \(k=1\), \(P_\varnothing(x)=\omega^x\), and
\(C=-r(\varnothing)\). 
Let \(E_\varnothing=\R_p\), and define its
\(p\) children by
\[
 E_a=\{x\in\R_p:\omega^x=\zeta_a\}=\{a\},
 \qquad a=0,\ldots,p-1,
\]
since \(\zeta_a=\omega^a\) and \(\omega\) is a primitive \(p\)-th
root of unity. These singletons partition \(E_\varnothing\), each
with normalized measure \(1/p\), and
\(P_\varnothing=\sum_a\zeta_a1_{E_a}=\omega^x\).
The only nonzero Fourier coefficient of \(P_\varnothing\) has modulus one.
All assertions hold for the root and its children; in particular,
the root exponent is \(C+r(\varnothing)=0\).
This completes round \(0\): the depth-zero polynomial and the sets
through depth one have been constructed.

\emph{Round \(j\), where \(1\leq j<N\).}
Suppose rounds \(0,\ldots,j-1\) have been completed on \(\R_p^n\).
Thus the sets \(E_v\) are available for \(|v|\leq j\), the
polynomials \(P_v\) are available for \(|v|<j\), and properties
\emph{(1)}--\emph{(4)} hold for all objects constructed so far, with
a common function \(k\) and integer \(C\). Round \(j\) constructs
simultaneously the polynomials \(P_v\), \(|v|=j\), and the sets
\(E_{va}\) associated with the child nodes \(va\), \(|va|=j+1\).

The round has three operations. First, a common chirp shifts all
previous coefficient levels by the same amount. Second, a localized
chirp realizes each depth-\(j\) node on its prescribed set \(E_v\).
Finally, a private coordinate separates its spectrum from every other
node and divides \(E_v\) into \(p\) children of equal measure.

All existing polynomials have coefficient levels
\(p^{-(C+r(v))}\). Choose an integer \(M\geq0\) large enough that
\begin{equation}
 d_v:=C+M+r(v)\geq n\qquad (|v|=j).
\end{equation}
For example, take
\[
 M=\max\bigl\{0,\ n-C-\min_{|v|=j}r(v)\bigr\}.
\]
On \(2M\) fresh coordinates introduce the unimodular chirp
\(Q(u,w)=\omega^{u\cdot w}\). Replace
\[
 k\ \text{by}\ k'=Qk,\qquad
 P_v\ \text{by}\ P'_v=QP_v,\quad |v|<j,
 \qquad C\ \text{by}\ C'=C+M.
\]
The old identities \(P'_v=k'h_v\) hold with the \emph{same}
functions \(h_v\) and the same sets. By Lemma~\ref{lem:chirp} with
\(E=\R_p^M\), every coefficient of \(Q\) has modulus \(p^{-M}\).
Since the old variables and the variables of \(Q\) are disjoint,
the Fourier coefficients of \(QP_v\) are tensor products of those
of \(Q\) and \(P_v\). Thus every old coefficient level is multiplied
by \(p^{-M}\). Tensoring each old spectrum with the same set of new
frequencies also preserves pairwise disjointness.
Since \(M\geq0\), no old exponent decreases; each new exponent
\(C'+r(v)=d_v\) is at least \(n\geq1\).
This verifies nonnegativity of all exponents at every round.

For \(|v|<j\), \(P'_v\) denotes the leveled-up polynomial \(QP_v\)
just defined. We now reuse the same primed notation for each new
depth-\(j\) node, but with a different defining formula, since no
old polynomial exists yet at such a node; both cases are unprimed
together at the end of the round.

For each new node \(v\) of depth \(j\), introduce its own fresh
coordinates
\[
 z_v\in\R_p^{d_v-n},\qquad
 Y_v\in\R_p^{d_v},\qquad s_v\in\R_p.
\]
All these coordinate blocks are disjoint, both from one another
and from those of \(Q\). Let \(x\in\R_p^n\) denote the entire
pre-round ambient variable, and set \(X_v=(x,z_v)\). Thus \(x\)
contains only the coordinates present before the common-chirp extension,
so \(X_v\) does not include the newly introduced coordinates \((u,w)\).
 We regard \(E_v\)
as a cylinder set in the enlarged product group and define
\begin{equation}
 P'_v=1_{E_v}(x)\omega^{X_v\cdot Y_v+s_v}.
\end{equation}
Consequently, \(P'_v\) is independent of \((u,w)\).
Ignoring the last character for a moment, Lemma~\ref{lem:chirp}
gives coefficient magnitudes \(0\) or \(p^{-d_v}\). Multiplication by
\(\omega^{s_v}\) merely shifts the frequency in the private coordinate
\(s_v\) from zero to one, without changing these magnitudes. Hence
\(P'_v\) has exactly the desired coefficient level
\(p^{-d_v}=p^{-(C'+r(v))}\). Its independence of all other newly
introduced coordinate blocks forces frequency zero in those blocks.


The coordinate \(s_v\) verifies spectral disjointness without any
assumption on the shape of \(E_v\). Every nonzero coefficient of
\(P'_v\) has \(s_v\)-frequency one. All old polynomials and all
other polynomials introduced in this round are independent of
\(s_v\), and therefore
have frequency zero there. Consequently the new spectra are
pairwise disjoint and are disjoint from every old spectrum.
In later rounds, tensoring existing polynomials with a common chirp
on fresh coordinates preserves every separation already established.
A subsequently introduced node \(t\) is separated from all existing
nodes by its own fresh coordinate \(s_t\); it need not have
frequency zero in the earlier coordinate \(s_v\).
The same observation proves \(\widehat P'_v(0)=0\).
The \(s_v\)-frequency of an existing nonroot polynomial remains one
under every later multiplication by a fresh-coordinate chirp.
For the root, the original \(x\)-frequency similarly remains one.
Thus the zero-mean property also persists throughout the construction.

It remains to define the new children. 
We regard the previously constructed set \(E_v\) and function \(k'\) on the fully enlarged product group by cylindrical extension. On \(E_v\), put
\[
 h'_v:=\overline{\bigl(k'\bigr)}\,P'_v,
\]
where the bar denotes complex conjugation.
Define
\[
 E_{va}=\{x_{\rm full}\in E_v:h'_v(x_{\rm full})=\zeta_a\}.
\]
Since \((k')^p=1\),
\(\overline{k'}=(k')^{-1}=(k')^{p-1}\), we obtain that \(P'_v=k'h'_v\).
Outside \(E_v\), put \(h'_v=0\).
On \(E_v\), the function \(h'_v\) is a \(p\)-th root of unity.
Conditioning on every coordinate except \(s_v\), it is a fixed
\(p\)-th root of unity times \(\omega^{s_v}\), since \(k'\) is
independent of \(s_v\). As \(s_v\) runs uniformly through \(\R_p\),
each \(p\)-th root occurs exactly once. Thus each child has exactly
\(1/p\) of its parent's measure. The children partition \(E_v\), and,
because \(h'_v=\zeta_a\) on \(E_{va}\),
\[
 h'_v=\sum_{a=0}^{p-1}\zeta_a1_{E_{va}}.
\]

This step uses
\[
 n'=n+2M+\sum_{|v|=j}(2d_v-n+1)
\]
coordinates in total: the common chirp uses \(2M\), while the node
\(v\) uses \((d_v-n)+d_v+1=2d_v-n+1\) private coordinates. This
number is finite.
Previously defined children are unchanged. Drop the primes from the
updated objects and regard them as living on \(\R_p^{n'}\). Properties
\emph{(1)}--\emph{(4)} now hold with the polynomials constructed
through depth \(j\) and the sets constructed through depth \(j+1\).
This completes round \(j\).

After round \(N-1\), the polynomials have been constructed for every
\(v\in\T_N\), and the sets \(E_v\) are available through depth \(N\).
Set \(D\) to be the value of \(n\) reached after round \(N-1\), so
that all constructed objects live on \(\R_p^D\), as claimed.
Under the successive cylinder extensions, the initial root set
\(E_\varnothing=\R_p\) has become
\(\R_p\times\R_p^{D-1}=\R_p^D\), in agreement with property~\emph{(1)}.
\end{proof}

\end{document}
